\chapter{Összegzés}

\section{Mit mutat be a dolgozat?}

A dolgozat egy olyan rendszert mutat be, amely valós League of Legends mérkőzésadatokból indul ki,
majd ezekből játékosprofilokat, kapcsolatokat, gráfokat, klasztereket, elemzési eredményeket és
végül egy szimulációs kísérletet épít. Közérthetően fogalmazva: a rendszer nem csak azt nézi meg,
hogy egy játékos mennyit öl vagy hányszor hal meg, hanem azt is, hogy kikkel játszik együtt,
kikkel kerül szembe, milyen ismétlődő kapcsolatok rajzolódnak ki körülötte, és ezekből milyen
nagyobb hálózati minták olvashatók ki.

Ha valamit le kellene szűkíteni a projekt lényegéből, az inkább az egész lánc megléte, mint bármelyik egyedi komponens. Önmagában egy Dijkstra-implementáció vagy egy React-oldal nem különösebben érdekes. De ha a nyers meccsadatoktól el lehet jutni gráfon, klasztereken és Rust-analitikán át egy reprodukálható szimulációs kísérletig -- és minden lépés ellenőrizhető és újrafuttatható --, akkor az már más jellegű eredmény.

\section{A rendszer közérthető képe}

A League of Legends mérkőzésekben minden játékos egy nagyobb kapcsolati tér részévé válik.
Egyes játékosok újra és újra egymás mellé kerülnek, mások rendszeresen egymás ellen játszanak.
Ezek az ismétlődések önmagukban nem árulnak el semmit a játékosok szándékáról vagy érzéseiről,
de alkalmasak arra, hogy kapcsolatként kezeljük őket egy gráfban, és arra, hogy ebből a
kapcsolathálóból valódi következtetéseket vonjunk le.

A projekt ebből a felismerésből indul. Nem egy egyszerű statisztikai oldalt épít, hanem egy
teljes feldolgozási láncot: a nyers API-adattól a normalizált játékosprofilokon, a gráfépítésen,
a klaszterezésen és a Rust-alapú gráfanalitikán át egészen egy evolúciós szimulációig, amelybe
ezek a valódi adatból levezetett profilok épülnek be.

A kérdések, amelyeket a rendszer vizsgál, nem pszichológiai állítások. Vannak-e sűrűbb csoportok
a kapcsolathálóban? Vannak-e játékosok, akik sok csoport között közvetítenek, és ezek eltávolítása
szétszakítaná a hálót? Hasonló teljesítményű játékosok kerülnek-e egymás mellé szívesebben?
Másképpen viselkedik-e egy Flex Queue-ból felépített gráf (ahol az ismétlődő együttjátszás
természetes), mint egy SoloQ kontroll-gráf, ahol a matchmaking véletlenszerűsége dominál?

Ezek mérhető, módszertanilag védett kérdések. A válaszok nem mutatnak be valódi emberi szándékot,
de megmutatják, milyen hálózati szerkezetek keletkeznek adott adatgyűjtési és elemzési szabályok
mellett. Ez a különbségtétel (mit mér a rendszer és mit nem) az egész projekt tudományos
értékének alapja.

\section{A fő eredmények}

A dolgozat első nagy eredménye a több-adatkészletes adatgyűjtési és feldolgozási architektúra.
A \textit{flexset} és a \textit{soloq} adatkészlet külön kezelése azért fontos, mert a két queue
nem ugyanazt a jelenséget méri. A Flex Queue alkalmasabb ismétlődő kapcsolati minták vizsgálatára,
míg a SoloQ inkább kontrollként és egyéni teljesítményprofilként értelmezhető.

A második fontos eredmény az \textit{opscore} és \textit{feedscore} metrikák használata.
Ezek nem tökéletes igazságmérők, hanem értelmezhető, szerepkörérzékeny pontszámok. Az értékük abban
van, hogy a játékosok teljesítménye nem egyetlen nyers statisztikára szűkül, hanem több játékbeli
szempontból válik összehasonlíthatóvá. Ez tette lehetővé, hogy a teljesítménymutatók később
gráfanalitikai bemenetként is használhatók legyenek.

A harmadik eredmény a játékoskapcsolati gráf és annak elemzése. A \textit{flexset} adatkészleten
kirajzolódó szerkezet core-periphery jellegű: vannak sűrűbb, erősebben összekapcsolt magrészek,
és sok lazább, kisebb periférikus csoport. Ez a megfigyelés nem azt jelenti, hogy a gráf teljes
társas valóságot ír le, hanem azt, hogy az ismétlődő mérkőzéskapcsolatok nem véletlenül szétszórt
pontokként jelennek meg.

A negyedik eredmény az asszortativitási elemzés. Az \textit{opscore} esetében pozitív kapcsolat
figyelhető meg: a kapcsolt játékosok több esetben hasonlóbb teljesítményprofilt mutatnak, mint amit
teljesen véletlen kapcsolatoknál várnánk. A \textit{feedscore} viszont érzékenyebb a gráfmódra,
ami jól mutatja, hogy az elemzésben a projekciós döntések nem technikai apróságok, hanem az
eredmény értelmezését is alakítják.

Az ötödik eredmény a Rust-alapú algoritmikus réteg. A rendszer nem csak adatokat jelenít meg,
hanem gráfalgoritmusokat is futtat: útvonalkeresést, asszortativitást és betweenness-központiságot.
A projekt ezzel nem marad meg dashboard-szinten.
A számításigényes részek külön futtatókörnyezetben, determinisztikus és újrafuttatható módon
jelennek meg.

\section{A Tribal NeuroSim szerepe}

A dolgozat záró nagy egysége a Tribal NeuroSim. Ennek szerepe nem az, hogy közvetlenül
megjósolja a League of Legends meccsek kimenetelét. A NeuroSim inkább egy transfer-kísérlet:
azt vizsgálja, hogy a gráfból és a klaszterekből kinyert teljesítményprofilok átvihetők-e egy
evolúciós, többügynökös szimuláció kezdeti paramétereibe.

Egyszerűbben: a PremadeGraph megmondja, milyen profilú csoportokat látunk a játékosgráfban,
a NeuroSim pedig ezekből a profilokból induló törzseket enged versenyezni egy kontrollált
szimulációs világban. A törzsek nem valódi játékosokat jelentenek, hanem klaszterekből képzett
modell-entitásokat. Ez a különbség lényeges, mert így a szimuláció nem pszichológiai állítás,
hanem ellenőrizhető modellkísérlet marad.

A négy validált futtatás alapján a Flex Queue-ból származó profilok konvergensebb viselkedést
mutattak: mindkét vizsgált seed mellett hasonló, Warband/Combat jellegű győztes archetípus jelent
meg. A SoloQ futtatások változékonyabbak voltak: eltérő seedek mellett eltérő győztes profilok
alakultak ki. Ez nem bizonyítja, hogy a valódi Flex Queue játékosok ilyenek, vagy hogy a SoloQ
játékosok más pszichológiai típusba tartoznának. Azt viszont alátámasztja, hogy a két adatkészletből
származó klaszterpriorok ugyanabban a szimulációs szabályrendszerben eltérő emergáló dinamikához
vezethetnek.

\section{Módszertani határok}

A dolgozat egyik fontos erőssége, hogy nem próbál többet állítani annál, amit az adatok és a
módszerek elbírnak. A signed balance, vagyis az előjeles gráfok strukturális egyensúlya például
technikailag érdekes kísérleti irány marad, de nem fő empirikus eredmény. Ennek oka, hogy az
ellenségként modellált mérkőzéskapcsolat nem feltétlenül jelent valódi negatív társas kapcsolatot.
Lehet matchmaking-hatás, queue-időzítés, rangsorbeli közelség vagy adatgyűjtési következmény is.

Ugyanezért nem szerepelnek a dolgozat fő állításai között időbeli stabilitási vagy közösségi
kohézióra vonatkozó erős kijelentések. A projekt végső formája inkább a védhetőbb kérdésekre
koncentrál: hogyan épül a gráf, milyen teljesítménymetrikák kapcsolhatók hozzá, milyen
asszortativitási és központisági minták jelennek meg, hogyan különbözik a Flex Queue és a SoloQ,
és hogyan használhatók a klaszterprofilok egy reprodukálható szimuláció bemeneteként.

\section{Végső következtetés}

A dolgozat végső állítása az, hogy kompetitív játékadatokból építhető olyan rendszer, amely egyszerre
szoftvermérnöki és kutatási szempontból is értelmezhető. A PremadeGraph nem pusztán statisztikai
oldal, hanem adatgyűjtési, adatmodellezési, gráfanalitikai és szimulációs keretrendszer.

Ami a rendszerből visszatekintve a legjobban működött, az a lépések közötti átjárhatóság: a nyers meccsadatokból normalizált játékosprofil lesz, abból kapcsolati gráf, a gráfból klaszterek és analitikai mutatók, a klaszterekből pedig szimulációs priorok. Nem minden ilyen összekapcsolás volt eleve tervezett -- egy részük abból következett, hogy az előző lépés eredménye logikus bemenet volt a következőhöz.

Az, amit ez a projekt megmutat, talán nem is kimondottan a League of Legends-ről szól. Inkább arról, hogy strukturált interakciós adatokból -- ha megfelelő módszertannal közelítünk hozzájuk -- olyan hálózati elemzés építhető, ami nem pusztán leíró, hanem kérdéseket is fel tud tenni. Hogy ezek a kérdések pontosan mit fednek fel a játékosokról, az más kérdés -- és a dolgozat nem is próbál ennél messzebbre menni.

\section{Túl a videojátékon: módszertani transzferálhatóság}
\label{sec:osszegzes-transfer}

A PremadeGraph kiindulópontja egy League of Legends API-n keresztül gyűjtött adatbázis. Ez
tudatos választás volt, nem véletlenszerű: a Riot Games Developer API nyilvánosan elérhető,
gazdagon dokumentált, fejlesztői szinten kezelhető, és strukturált JSON adatot szolgáltat.
A játékadatok jellegük miatt ideálisan alkalmasak gráf-pipeline prototipizálásra: sok szereplő,
sok interakció, jól definiált szerepkörök, kvantitatív teljesítménymutatók. A módszertan
nem LoL-specifikus. Amit ez a projekt épített, az egy általános célú kutatási infrastruktúra,
amelynek minden komponense transzferálható.

\textbf{Gráfklaszterezés és közösségi struktúra.} A greedy modularity maximization ugyanolyan
természetességgel alkalmazható tudományos kollaborációs hálózatokra (kik dolgoznak együtt
rendszeresen?), vállalati szervezeti gráfokra (melyik csapatok alkotnak valódi munkaegységet
a formális orgchart mögött?), logisztikai hub-hálózatokra (mely elosztóközpontok alkotnak
regionális klasztert?) vagy közösségi média interakciókra (milyen buborékok rajzolódnak ki?).

\textbf{Asszortativitási elemzés.} Annak mérése, hogy hasonló attribútumú csomópontok
hajlamosabbak-e egymással kapcsolódni, általánosan alkalmazható. Pénzügyi hálózatokban
a kockázati kontagion terjedési mintáját írja le; ökológiai hálózatokban a niche-specializáció
fokát méri; epidemiológiai modellekben a kontaktstruktúra homofíliáját jellemzi.

\textbf{Párhuzamos Brandes-féle betweenness centralitás.} A kritikus közvetítő csomópontok
azonosítása infrastruktúra-tervezésben (melyik útcsomópont kiesése bénítja meg a hálózatot?),
járványvédelemben (mely egyén izolálása csökkenti leginkább a terjedést?) és
ellátásilánc-elemzésben (melyik raktár kiesése okoz maximális fennakadást?) egyaránt
közvetlenül alkalmazható.

\textbf{Pathfinder algoritmusok.} A BFS, Dijkstra, kétirányú keresés és az egzakt A* nem
League of Legends találmányok. Navigációs rendszerek, hálózati routing optimalizáció,
robotikai útvonaltervezés, játékmotorok: mindegyik ugyanezt az algoritmikus rétegt
használja. A Pathfinder Lab összehasonlítási felülete adaptálható bármely gráf-alapú
keresési problémára.

\textbf{3D Bird's Eye Sphere és Fibonacci-vetítés.} A gömbfelületen egyenletesen elosztott
gráf-vizualizáció egyedi megközelítést kínál nagy hálózatok struktúrájának áttekintéséhez.
Ahol a klasszikus 2D force-directed elrendezés átfedő, olvashatatlan ábrákat produkál,
a 3D gömbvetítés fenntartja a csúcsok szétválaszthatóságát. Bioinformatikában fehérje-interakciós
hálózatoknál, szociológiában globális kapcsolathálózatoknál ez kifejezetten hasznos eszköz.

\textbf{Tribal NeuroSim evolúciós szimulációja.} Az artifact-prior alapú inicializálás
és a NEAT-stílusú neuroevolúció kombinációja bármilyen domainre alkalmazható, ahol
csoportszintű teljesítményprofilok kinyerhetők és összehasonlíthatók. Az öt
alapartifakt (combat, resource, map\_objective, risk, team) nem LoL-fogalom.
Csapatsportokban harci~=~fragtörténet, resource~=~gazdasági döntések, map~=~pozícionálás,
risk~=~agresszív húzások vállalása, team~=~asszisztok és csapatmentések. Startupok esetén
combat~=~piaci agresszivitás, resource~=~tőkehatékonyság, map~=~piaci célzás, risk~=~kockázattűrés,
team~=~csapatkoordináció. Az ökológiában combat~=~predációs képesség, resource~=~táplálékkeresési
hatékonyság, team~=~koloniális koordináció. A rendszer ezeket az analógiákat nem spekulálja --
hanem megmutatja az architektúrát, amellyel bármely ilyen adatforrás ugyanerre a szimulációs
pipeline-ra illeszthető.

A Riot Games API azért volt kényelmes prototípus-forrás, mert nyílt, strukturált és gazdag
metaadatokban. De a módszertan nem igényel LoL-adatot. Ugyanez a pipeline futtatható bármely
nyílt API-ra (GitHub kollaborációs gráf, NBA teljesítményadatbázis, PubMed szerzői hálózat),
ahol hasonló felosztás elvégezhető. A videojáték itt eszköz volt, nem korlát.

Ahogy az ABBA megfogalmazta évtizedekkel ezelőtt: \textit{the winner takes it all}.
A Tribal NeuroSimben ez nem metafora, hanem a mechanika közvetlen következménye, és egyben
az egész pipeline szimbóluma. Az adat, a gráf, a klaszter, az analitika és a szimuláció mind
egyetlen coherens lánc részei. Aki megérti ezt a láncot, annak kezében van egy általánosan
alkalmazható hálózatelemzési és evolúciós modellezési keretrendszer, amelynek a videojáték
csak a legpraktikusabb és legelérhetőbb bemeneti forrása volt.
